% !TeX root = ../main.tex
\section*{Заключение}

В ходе выполнения работы была разработана и исследована имитационная модель тракта цифровой демодуляции и согласованной фильтрации радиосигналов в среде MATLAB для параметров Варианта~8 ($f_c = 210.0$~МГц, $f_s = 180.0$~МГц, $n_a = n_g = 13$~бит). Сформированный тракт обеспечил корректный квадратурный перенос ЛЧМ-сигнала на нулевую промежуточную частоту и подавление внеполосных составляющих КИХ-фильтром нижних частот.

Моделирование подтвердило строгую математическую эквивалентность прямой свертки во временной области (FIR), спектрального метода (FFT/IFFT) и блочной сегментированной свертки методом перекрытия с накоплением (\texttt{fftfilt}), графики откликов которых полностью совпали. При этом метод \texttt{fftfilt} является наиболее предпочтительным для аппаратной реализации в реальном времени, поскольку сочетает вычислительную эффективность алгоритмов БПФ с возможностью потоковой обработки без задержки на накопление всей выборки.

Исследование оконной весовой обработки показало, что применение окна Хэмминга позволяет снизить уровень боковых лепестков с теоретических $-13.3$~дБ до уровня менее $-42.8$~дБ, что предотвращает маскирование слабых эхо-сигналов ценой контролируемого расширения главного пика сжатия в $1.35$--$1.4$ раза.
